\documentclass[11pt]{article}
\usepackage{mathblatt}
% Prüfstein M2 09.10.2026, Lernblatt: „Kathete berechnen“ (Pythagoras, Lerneinheit 2, erster Teil;
% die Umkehrung ist ein eigenes Blatt). Nach bau/bauauftrag.md von oben gebaut; Muster K4W.
% Lernweg: L2k-1 Idee (Flächen abziehen) · L2k-2 Hypotenuse oder Kathete gesucht · L2k-3 Gleichung
% umstellen ohne Zahl, dann mit Wurzel wählen · L2k-4 Rechnen glatt, a) grau, Probe · L2k-5 Rechenaufwand
% kurz · L2k-6 Sache: Bild ohne Dreieck → Karte → nur Text · L2k-7 Ziel nach P10 ’24 (Seilbahn).
% Gemeinsame Präambel, übernommen aus hypotenuse-berechnen (K4W, Probe 08.10.2026).
\makeatletter
\setlength{\mbpfbe}{0mm}% keine Punkte (Bauregel 6.4)
% eingerückt (Bauregel 4.7, 6.6): \gEin Einzug, \gSchrift eine Stufe kleiner
\newlength{\gEin}\setlength{\gEin}{0pt}
\let\gSchrift\relax
\renewcommand{\pfkreuzzeile}[1]{\par\noindent\hangindent3.2em\hangafter1\hspace*{1.6em}$\square$\ #1\par}
\newcommand{\gkreuz}[1]{$\square$\ #1\hspace{2em}}
% Bauregel 6.7: links, was man ansieht, rechts, was man tut; Spaltengrenze überall an derselben Stelle
\newsavebox{\gbox}\newlength{\gLinks}
\newcommand{\gzwei}[2]{\par\smallskip\noindent\sbox{\gbox}{#1}%
  \setlength{\gLinks}{\dimexpr0.5\textwidth-\gEin-\mbpfnr\relax}%
  \ifdim\wd\gbox>\dimexpr\gLinks-4mm\relax
    \usebox{\gbox}\par\smallskip\noindent\begin{minipage}{\linewidth}#2\end{minipage}%
  \else
    \begin{minipage}[t]{\gLinks}\vspace{0pt}\usebox{\gbox}\end{minipage}%
    \begin{minipage}[t]{\dimexpr\linewidth-\gLinks\relax}\vspace{0pt}\raggedright #2\end{minipage}%
  \fi\par}
\RenewDocumentEnvironment{pfaufg}{m m m}{%
  \begin{lrbox}{\mbaufgabebox}\begin{minipage}[t]{\linewidth}%
  \gSchrift\noindent\hspace*{\gEin}\mbpfrandmarke{#1}%
  \begin{minipage}[t]{\mbpfnr}\strut\textbf{#2}\end{minipage}%
  \begin{minipage}[t]{\dimexpr\linewidth-\gEin-\mbpfnr\relax}\raggedright\strut\ignorespaces}%
 {\end{minipage}%
  \end{minipage}\end{lrbox}%
  \setlength{\mbaufgabehoehe}{\dimexpr\ht\mbaufgabebox+\dp\mbaufgabebox\relax}%
  \par\addvspace{7pt}\Needspace*{\mbaufgabehoehe}%
  \noindent\usebox{\mbaufgabebox}\par\addvspace{7pt}}
\makeatother
% Rechenraum als Karo (Bauregel 6.9): n Rechenschritte = n mal 10 mm
\newcommand{\karo}[1]{\par\smallskip\noindent\begin{tikzpicture}
  \clip (0,0) rectangle ({\linewidth-1mm},{#1*10mm});
  \draw[black!16,line width=0.3pt,step=5mm] (0,0) grid ({\linewidth},{#1*10mm});
  \end{tikzpicture}\par}
\newcommand{\kk}[1]{$\square$\,#1\hspace{0.9em}}
\newcommand{\linie}[1]{\,{\color{black!45}\rule[-1pt]{#1}{0.4pt}}\,}
% Zeile am Ende jedes Durchgangs: Originale klein grau, darüber; dann Vorgänger/Nachfolger (Bauregel 3.4, 6.5)
\newcommand{\blattende}[2]{\par\vfill\noindent\begin{minipage}{\linewidth}{\scriptsize\color{mbgrau}Originale: #1\par}\smallskip
{\footnotesize #2\par}\end{minipage}}
\newcommand{\pkt}{\textperiodcentered{}}

\newcommand{\kennung}{T6B}
\newcommand{\grau}[1]{{\color{mbgrau}#1}}

\begin{document}
\pfheftstil
\fancyfoot[C]{\parbox[t]{\textwidth}{\mbfhbox\par\vspace{1.5mm}{\footnotesize\color{mbgrau}{\scriptsize \kennung}\hfill\thepage}}}
\pfheftkopf{Kathete berechnen}{Klasse 9}

% L2k-1 Idee: der Satz rückwärts über Flächen – vom großen Quadrat abziehen (Fehler: Quadrate addiert)
\begin{pfaufg}{}{1.}{}
Die zwei kleinen Quadrate über den Katheten haben zusammen so viel Fläche wie das große Quadrat über der Hypotenuse. Hier fehlt eine Fläche.
\gzwei{\begin{tikzpicture}[scale=0.85,line width=0.7pt]
\draw (0,0) -- (2,0) -- (2,-2) -- (0,-2) -- cycle;
\draw (0,0) -- (-1.5,0) -- (-1.5,1.5) -- (0,1.5) -- cycle;
\draw (0,1.5) -- (2,0) -- (3.5,2) -- (1.5,3.5) -- cycle;
\draw[line width=1.4pt] (0,0) -- (2,0) -- (0,1.5) -- cycle;
\mbrwmarke{(0,0)}{(2,0)}{(0,1.5)}
\node[font=\small] at (1.75,1.75) {$100$ cm$^2$};
\node[font=\footnotesize,align=center] at (-0.75,0.75) {$36$\\cm$^2$};
\node[font=\small] at (1,-1) {?};
\end{tikzpicture}}{%
\textbf{a)}\ Wie groß ist die Fläche des Quadrats mit dem Fragezeichen?\linie{14mm}cm$^2$\par\bigskip
\textbf{b)}\ Wie lang ist die Seite dieses Quadrats?\linie{14mm}cm}
\fusshilfe{\mbox{1:~a) 64 cm²; b) 8 cm}}
\end{pfaufg}

% L2k-2 Erkennen ohne Rechnen: gemischt, verschiedene Lagen (Fehler: längere Kathete für die Hypotenuse halten)
\begin{pfaufg}{}{2.}{}
Zwei Seiten sind gegeben. Ist die Seite mit dem Fragezeichen die Hypotenuse oder eine Kathete? Kreuze an.\par\medskip\noindent
\begin{minipage}[b]{0.245\linewidth}\centering
\begin{tikzpicture}[line width=0.7pt,scale=0.85]\coordinate (P1) at (0,0);\coordinate (P2) at (2.0,0);\coordinate (P3) at (0,1.5);\draw (P1) -- (P2) -- (P3) -- cycle;\mbrwmarke{(P1)}{(P2)}{(P3)}
\node[font=\small,anchor=north] at (1.0,-0.08) {$8$ cm};\node[font=\small,anchor=east] at (-0.08,0.75) {$6$ cm};\node[font=\small,anchor=south west] at (1.02,0.78) {?};\end{tikzpicture}\par\smallskip
{\small a)\ \kk{Hypotenuse}\par\hspace*{1.1em}\kk{Kathete}}
\end{minipage}\hfill
\begin{minipage}[b]{0.245\linewidth}\centering
\begin{tikzpicture}[line width=0.7pt,scale=0.85]\coordinate (P1) at (0,0);\coordinate (P2) at (2.55,0);\coordinate (P3) at (1.985,1.059);\draw (P1) -- (P2) -- (P3) -- cycle;\mbrwmarke{(P3)}{(P1)}{(P2)}
\node[font=\small,anchor=north] at (1.27,-0.08) {$17$ cm};\node[font=\small,anchor=south east] at (0.95,0.56) {$15$ cm};\node[font=\small,anchor=south west] at (2.30,0.56) {?};\end{tikzpicture}\par\smallskip
{\small b)\ \kk{Hypotenuse}\par\hspace*{1.1em}\kk{Kathete}}
\end{minipage}\hfill
\begin{minipage}[b]{0.245\linewidth}\centering
\begin{tikzpicture}[line width=0.7pt,scale=0.85]\coordinate (P1) at (0,0);\coordinate (P2) at (2.4,0);\coordinate (P3) at (2.4,0.7);\draw (P1) -- (P2) -- (P3) -- cycle;\mbrwmarke{(P2)}{(P1)}{(P3)}
\node[font=\small,anchor=north] at (1.2,-0.08) {?};\node[font=\small,anchor=west] at (2.48,0.35) {$7$ cm};\node[font=\small,anchor=south east] at (1.18,0.38) {$25$ cm};\end{tikzpicture}\par\smallskip
{\small c)\ \kk{Hypotenuse}\par\hspace*{1.1em}\kk{Kathete}}
\end{minipage}\hfill
\begin{minipage}[b]{0.245\linewidth}\centering
\begin{tikzpicture}[line width=0.7pt,scale=0.85]\coordinate (V) at (0.15,0.25);\coordinate (L) at (2.48,1.01);\coordinate (K) at (-0.12,1.07);\draw (V) -- (L) -- (K) -- cycle;\mbrwmarke{(V)}{(L)}{(K)}
\node[font=\small,anchor=north west] at (1.32,0.63) {$34$ m};\node[font=\small,anchor=east] at (-0.02,0.62) {$12$ m};\node[font=\small,anchor=south] at (1.18,1.06) {?};\end{tikzpicture}\par\smallskip
{\small d)\ \kk{Hypotenuse}\par\hspace*{1.1em}\kk{Kathete}}
\end{minipage}
\fusshilfe{\mbox{2:~a) H; b) K; c) K; d) H}}
\end{pfaufg}

% L2k-3 Gleichung umstellen ohne Zahl; a) grau vorgerechnet (Bauregel 1.5); wechselnde Buchstaben,
% Punktnamen wie in der P10 (Fehler: gesuchte Seite allein, „c ist immer die Hypotenuse“)
\begin{pfaufg}{}{3.}{}
Schreibe den Satz des Pythagoras für das Dreieck auf. Stelle die Gleichung dann so um, dass das Quadrat der gesuchten Seite allein steht.\par\medskip\noindent
\begin{minipage}[b]{0.32\linewidth}\centering
\begin{tikzpicture}[line width=0.7pt,scale=0.85]\coordinate (P1) at (0,0);\coordinate (P2) at (2.4,0);\coordinate (P3) at (0,1.5);\draw (P1) -- (P2) -- (P3) -- cycle;\mbrwmarke{(P1)}{(P2)}{(P3)}
\node[font=\small,anchor=north] at (1.2,-0.08) {$u$};\node[font=\small,anchor=east] at (-0.08,0.75) {$v$};\node[font=\small,anchor=south west] at (1.22,0.78) {$w$};\end{tikzpicture}\par\smallskip
\grau{{\small a)\ gesucht: $u$}\par\smallskip
$w^2 = u^2 + v^2$\par\smallskip $\Rightarrow\ u^2 = w^2 - v^2$}
\end{minipage}\hfill
\begin{minipage}[b]{0.32\linewidth}\centering
\begin{tikzpicture}[line width=0.7pt,scale=0.85]\coordinate (P1) at (0,0);\coordinate (P2) at (2.55,0);\coordinate (P3) at (0.565,1.059);\draw (P1) -- (P2) -- (P3) -- cycle;\mbrwmarke{(P3)}{(P1)}{(P2)}
\node[font=\small,anchor=north] at (1.27,-0.08) {$r$};\node[font=\small,anchor=south east] at (0.25,0.56) {$p$};\node[font=\small,anchor=south west] at (1.60,0.56) {$q$};\end{tikzpicture}\par\smallskip
{\small b)\ gesucht: $q$}\par\smallskip
\linie{30mm}\par\medskip\linie{30mm}
\end{minipage}\hfill
\begin{minipage}[b]{0.32\linewidth}\centering
\begin{tikzpicture}[line width=0.7pt,scale=0.85]\coordinate (A) at (0,1.5);\coordinate (B) at (0,0);\coordinate (C) at (2.4,0);\draw (A) -- (B) -- (C) -- cycle;\mbrwmarke{(B)}{(C)}{(A)}
\node[font=\small,anchor=south] at (A) {$A$};\node[font=\small,anchor=north east] at (B) {$B$};\node[font=\small,anchor=north west] at (C) {$C$};\end{tikzpicture}\par\smallskip
{\small c)\ gesucht: $\overline{BC}$}\par\smallskip
\linie{30mm}\par\medskip\linie{30mm}
\end{minipage}
\fusshilfe{\mbox{3:~b) $q^2 = r^2 - p^2$; c) $\overline{BC}^2 = \overline{AC}^2 - \overline{AB}^2$}}
\end{pfaufg}

% L2k-3, zweiter Teil: dieselbe Umstellung mit Wurzel, unter Fallen wählen (2024 · 1 · f ganz)
\begin{pfaufg}{P10 ’24}{4.}{}
Im Dreieck sind $y$ und $z$ die Katheten, $x$ ist die Hypotenuse. Kreuze die Gleichung an, mit der man die Länge von $z$ berechnen kann.
\gzwei{\begin{tikzpicture}[line width=0.7pt,scale=0.85]\coordinate (P1) at (0,0);\coordinate (P2) at (2.6,0);\coordinate (P3) at (0,1.6);\draw (P1) -- (P2) -- (P3) -- cycle;\mbrwmarke{(P1)}{(P2)}{(P3)}
\node[font=\small,anchor=north] at (1.3,-0.08) {$z$};\node[font=\small,anchor=east] at (-0.08,0.8) {$y$};\node[font=\small,anchor=south west] at (1.32,0.83) {$x$};\end{tikzpicture}}{%
\kk{$z = x^2 - y^2$}\kk{$z = \sqrt{x^2 + y^2}$}\par\medskip
\kk{$z = x + y$}\kk{$z = \sqrt{x^2 - y^2}$}}
\fusshilfe{\mbox{4:~$z = \sqrt{x^2 - y^2}$}}
\end{pfaufg}

% L2k-4 Rechnen mit glatten Zahlen; a) grau vorgerechnet, b) c) selbst; Probe kürzer als die Hypotenuse
\begin{pfaufg}{}{5.}{}
Die Hypotenuse und eine Kathete eines rechtwinkligen Dreiecks sind gegeben. Berechne die andere Kathete.\par\medskip
\grau{a)\ Hypotenuse $10$ cm, Kathete $6$ cm:\quad $a^2 = 10^2 - 6^2 = 100 - 36 = 64$\quad $\Rightarrow$\quad $a = \sqrt{64} = 8$ cm}\par\medskip
\noindent\begin{minipage}[t]{0.48\linewidth}
b)\ Hypotenuse $13$ cm, Kathete $5$ cm
\karo{3}
\end{minipage}\hfill
\begin{minipage}[t]{0.48\linewidth}
c)\ Hypotenuse $25$ cm, Kathete $7$ cm
\karo{3}
\end{minipage}
\par Prüfe dein Ergebnis: Die Kathete ist kürzer als die Hypotenuse.
\fusshilfe{\mbox{5:~b) 12 cm; c) 24 cm}}
\end{pfaufg}

% L2k-5 Rechenaufwand erst nach allen Denkschritten, kurz: Wurzel geht nicht auf; Einheiten angleichen
\begin{pfaufg}{}{6.}{}
Berechne die fehlende Kathete. Runde, wo nötig, auf eine Stelle nach dem Komma.\par\smallskip
\noindent\begin{minipage}[t]{0.48\linewidth}
a)\ Hypotenuse $8$ cm, Kathete $3$ cm
\karo{3}
\end{minipage}\hfill
\begin{minipage}[t]{0.48\linewidth}
b)\ Hypotenuse $1{,}3$ m, Kathete $50$ cm
\karo{3}
\end{minipage}
\fusshilfe{\mbox{6:~a) 7,4 cm; b) 120 cm}}
\end{pfaufg}

% L2k-6 Sache: der Schüler zeichnet das Dreieck selbst – erst im Bild ohne Dreieck (Seil nicht gezeichnet),
% dann auf einer Karte (Punkt gegenüber selbst finden), dann nur Text; Fragen wechseln: reicht es? (nein,
% K4W hatte ja) · Hypotenuse gesucht, Karte (Kritiker: eine gemischte Aufgabe) · wie hoch, mit Zusatzstrecke
\begin{pfaufg}{}{7.}{}
Tim spannt ein Seil von der Spitze der Zeltstange bis zu einem Hering im Boden. Das Seil ist $3{,}4$ m lang. Reicht der Platz bis zum Zaun dafür?
\gzwei{\begin{tikzpicture}[line width=0.7pt,scale=1.0]
\draw (-1.4,0) -- (3.6,0);
\draw[line width=1.8pt] (0,0) -- (0,1.6);
\foreach \x in {2.8,3.05,3.3,3.55} {\draw (\x,0) -- (\x,0.9);}
\draw (2.7,0.75) -- (3.6,0.75);\draw (2.7,0.3) -- (3.6,0.3);
\node[font=\small,anchor=east] at (-0.05,1.85) {Stange};
\draw[|-|,line width=0.4pt] (0,-0.35) -- (2.8,-0.35) node[midway,below,font=\small] {$2{,}8$ m};
\draw[|-|,line width=0.4pt] (0.35,0) -- (0.35,1.6) node[midway,right,font=\small] {$1{,}6$ m};
\node[font=\small,anchor=south] at (3.15,0.95) {Zaun};
\end{tikzpicture}}{%
\karo{3}\par\smallskip\hfill\kk{ja}\kk{nein}}
\fusshilfe{\mbox{7:~nein, gebraucht werden 3 m}}
\end{pfaufg}

\begin{pfaufg}{}{8.}{}
Ein Boot fährt bei $A$ los und will genau gegenüber anlegen. Die Strömung treibt es ab. Es legt bei $B$ an, $90$ m flussabwärts. Wie lang war die Fahrt von $A$ nach $B$?
\gzwei{\begin{tikzpicture}[line width=0.7pt,scale=0.8]
\fill[black!10] (-0.3,0) rectangle (4.3,2.4);
\draw (-0.3,0) -- (4.3,0);\draw (-0.3,2.4) -- (4.3,2.4);
\draw[->,line width=0.5pt] (2.6,1.2) -- (3.6,1.2);\node[font=\scriptsize,anchor=south] at (3.1,1.2) {Strömung};
\fill (0.6,0) circle (1.8pt) node[font=\small,anchor=north] {$A$};
\fill (2.4,2.4) circle (1.8pt) node[font=\small,anchor=north west] {$B$};
\draw[|-|,line width=0.4pt] (0.6,2.75) -- (2.4,2.75) node[midway,above,font=\small] {$90$ m};
\draw[|-|,line width=0.4pt] (-0.65,0) -- (-0.65,2.4) node[midway,left,font=\small] {$120$ m};
\end{tikzpicture}}{%
\karo{3}}
\fusshilfe{\mbox{8:~150 m}}
\end{pfaufg}

\begin{pfaufg}{}{9.}{}
Lea lässt einen Drachen steigen. Die Schnur ist $65$ m lang und straff gespannt. Waagerecht gemessen ist der Drachen $39$ m von Lea entfernt. Ihre Hand ist $1$ m über dem Boden. Wie hoch über dem Boden fliegt der Drachen?\par
Skizziere zuerst die Lage.
\karo{4}
\fusshilfe{\mbox{9:~53 m}}
\end{pfaufg}

% L2k-7 Zielaufgabe (2024 · 6 · a, Punkt C und Teil b weggelassen): Höhe als Kathete, dreistellige Zahlen, Runden
\begin{pfaufg}{nach P10 ’24}{10.}{}
Eine Seilbahn fährt von der Talstation $B$ zur Bergstation $A$. Das Seil $\overline{AB}$ ist $384$ m lang. Der Punkt $F$ liegt senkrecht unter $A$. Die Strecke $\overline{FA}$ ist der Höhenunterschied zwischen $B$ und $A$. Berechne diesen Höhenunterschied. Runde auf eine Stelle nach dem Komma.
\gzwei{\begin{tikzpicture}[line width=0.7pt,scale=0.8]
\fill[black!8] (-1.0,0) -- (-0.4,2.2) -- (0,2.87) -- (0.5,2.3) -- (1.1,1.4) -- (1.5,0.6) -- (2.2,0) -- cycle;
\draw[black!45] (-1.0,0) -- (-0.4,2.2) -- (0,2.87) -- (0.5,2.3) -- (1.1,1.4) -- (1.5,0.6) -- (2.2,0);
\draw (-1.2,0) -- (4.3,0);
\coordinate (A) at (0,2.87);\coordinate (F) at (0,0);\coordinate (B) at (3.83,0);
\draw[dashed] (A) -- (F);\draw (F) -- (B);\draw[line width=1.2pt] (A) -- (B);
\mbrwmarke{(F)}{(B)}{(A)}
\node[font=\small,anchor=south] at (A) {$A$};\node[font=\small,anchor=north] at (F) {$F$};\node[font=\small,anchor=north] at (B) {$B$};
\node[font=\small,anchor=south west] at (2.0,1.45) {$384$ m};
\node[font=\small,anchor=north] at (1.92,-0.25) {$255$ m};
\end{tikzpicture}}{%
\karo{3}}
\fusshilfe{\mbox{10:~287,1 m}}
\end{pfaufg}

\blattende{2024 \pkt{} 1 \pkt{} f \quad 2024 \pkt{} 6 \pkt{} a}{Vorher: Hypotenuse berechnen\quad\textbullet\quad Weiter: Das Dreieck erst finden}
\end{document}
